\section{Introduction}\label{sec:intro}

Cultural institutions are among the most energy-intensive non-industrial
buildings. Climate-controlled museum spaces consume
250--500\,kWh\,m$^{-2}$\,yr$^{-1}$, two to three times the intensity of
comparable office space, because conservation envelopes for sensitive
collections require year-round narrow-band control of both temperature and
relative humidity \citep{ashrae_handbook_ch24,verticchio2024microclimate}.
In Europe alone, public-museum HVAC consumes an estimated 5\,TWh of
electricity annually, and double-carbon policy goals are now placing
pressure on cultural institutions to decarbonize without compromising
preventive conservation.

The standard institutional response is to operate at a single nominal
setpoint, typically $22^\circ\mathrm{C}$ and $50\%$ relative humidity, that
sits inside the conservation envelope of the strictest collection in the
building. This approach has three documented weaknesses. First, it ignores
time-varying loads: visitor-driven sensible and latent gains, solar gain,
and outdoor weather all vary by an order of magnitude across the day, so
a constant setpoint converts every disturbance into HVAC work. Second, it
underuses the per-zone heterogeneity of the collection. Standards such as
EN~15757~\citep{en15757}, the broader heritage-building specification
EN~16893~\citep{en16893}, and ASHRAE Class~A1/A2~\citep{ashrae_handbook_ch24}
prescribe ranges, not points, and the safe range for bronze
($\mathrm{RH}\in[30,40]\%$) is incompatible with the safe range for oil
paintings ($\mathrm{RH}\in[45,55]\%$); a single setpoint necessarily
violates one of them. Third, it treats conservation as a static envelope
and ignores the rate clauses in modern standards: EN~15757 limits hourly
$\mathrm{d}\mathrm{RH}/\mathrm{d}t$ for hygroscopic materials, because
mechanical damage is dominated by short-term cycles rather than long-term
mean states.

Reinforcement learning (RL) and model-predictive control (MPC) have been
studied extensively for office and residential HVAC
\citep{brandi2022online,wang2017drlhvac,drgona2020mpc_review}. Open
testbeds such as Sinergym~\citep{manjavacas2024sinergym} now make it
straightforward to plug EnergyPlus building models into Gymnasium-style
training loops. Deploying these controllers in museums, however, forces a
qualitative change in safety semantics: an office-comfort violation is
annoying, while a conservation-envelope violation accumulates irreversible
damage to irreplaceable objects. The museum-specific HVAC literature has
so far either used classical multi-objective optimization without learning
\citep{ascione2018museum,schito2020museum} or treated conservation limits
as soft penalties inside a standard RL reward
\citep{green_smart_museums_2025,boesgaard2022heritage}; the constrained
formulation remains underexplored.

\paragraph{Contributions.}
This paper makes five contributions, C1--C5. C1: museum HVAC control is
formalized as a Constrained
Markov Decision Process (CMDP) in which (i) static temperature and
humidity envelopes from ASHRAE Class~A1, A2, and unconstrained categories
and (ii) the EN~15757 dynamic rate clauses on
$|\mathrm{d}T/\mathrm{d}t|$ and $|\mathrm{d}\mathrm{RH}/\mathrm{d}t|$
are encoded as a scalar cost. To the best of our knowledge, this is the
first CMDP formulation that includes the rate clauses. C2: we implement
a Lagrangian-PPO controller with a PID-tuned multiplier following
\citet{stooke2020responsive}, and document a subtle but blocking
implementation issue: when training with subprocess-based parallel
environments, $\lambda$ must be propagated via the \texttt{VecEnv}
attribute-setter mechanism, not through a Python mutable reference
captured in an env-factory closure. C3: a five-zone differentiated museum
simulator is constructed with collection-specific conservation envelopes
and a synthetic visitor flow, and Lagrangian-PPO is benchmarked against
four baselines: rule-based fixed setpoint, schedule-relaxed rule-based,
vanilla unconstrained PPO, and fixed-weight soft-penalty PPO. C4: we
trace a three-objective Pareto front by sweeping the cost target. The
front shows that the major energy-conservation trade-off lies along the
conservation axis: relaxing the cost budget by an order of magnitude
changes annual energy by less than $1\%$ but more than doubles the
conservation cost. C5: we build a high-fidelity EnergyPlus testbed,
extending Sinergym's five-zone building with per-zone thermostats and
active humidification, and use it to stress-test the method. We report
honestly that the localization finding transfers but the performance
advantage does not: on the high-fidelity model the constrained
controller is outperformed by a tuned fixed setpoint, which we trace to
the energy gradient overwhelming the constraint signal. We release the
testbed and analyze what conditioning would be needed to close the gap.

The remainder of the paper is organized as follows.
Section~\ref{sec:related} positions the work against existing museum
HVAC optimization, Safe RL, and digital-twin building literatures.
Section~\ref{sec:problem} states the CMDP formulation and the
conservation cost. Section~\ref{sec:method} describes the Lagrangian-PPO
algorithm and the implementation issue noted in C2.
Section~\ref{sec:setup} describes the simulator and baselines.
Section~\ref{sec:results} reports the Pareto front and the per-zone
analysis. Section~\ref{sec:discussion} discusses limitations and the
path to a real-EnergyPlus deployment.

% =====================  2. RELATED WORK  ====================================
\section{Related Work}\label{sec:related}

\subsection{Museum HVAC optimization and microclimate}

Museum HVAC research has evolved from monitoring-only studies toward
control-aware multi-objective optimization. \citet{ascione2018museum}
formulated museum microclimate control as a three-objective problem in
energy, visitor discomfort, and artwork preservation risk, and solved it
with NSGA-II. Their preservation risk is computed offline from
monthly-aggregated environmental traces.
\citet{schito2020museum} extended this line of work with an
achievement-function formulation applied to a Pisa case study, producing
Pareto-optimal HVAC setpoint schedules. Both works treat the controller
as a parameter-tuning problem rather than a closed-loop policy, and they
do not include the EN~15757 rate clauses.
\citet{coelho2020museum} took a complementary route, using hygrothermal
simulation to assess the impact of climate change on artefact
conservation and building energy use, but operates offline. \citet{verticchio2024microclimate} re-examined
microclimate thresholds parametrically and quantified the energy cost of
each standard tolerance class.

Recent work has begun to combine learning with conservation-aware
control. \citet{boesgaard2022heritage} applied gradient-boosting and
LSTM models to predict indoor temperature and humidity evolution in two
heritage buildings. \citet{green_smart_museums_2025} proposed an
integrated framework that combines digital-twin modeling, LSTM-based load
forecasting, and a Deep Q-Network (DQN) controller for HVAC and lighting
across three Chinese museum case studies. Their DQN minimizes a weighted
sum of energy cost and visitor-comfort penalty; conservation enters only
as a soft penalty term and EN~15757 rate clauses are not formalized.
Outside the optimization stream, \citet{museum_no_hvac_2025},
\citet{greek_museum_2024}, and \citet{coherence_heritage_2025}
characterize microclimates in heritage buildings without active control,
while \citet{membrane_dehumid_museum_2025} explored hardware-side
interventions such as membrane-based dehumidification. The thermal
comfort side has also been re-examined for museums:
\citet{mudejar2024adaptive} showed that adaptive comfort models
\citep{ashrae552020,en167982020} can relax setpoint requirements without
increasing visitor dissatisfaction. None of these works formulate a
closed-loop constraint-respecting controller of the type studied here.

\subsection{Safe and constrained reinforcement learning}

Constrained Markov Decision Processes~\citep{altman1999constrained}
provide the canonical framework for RL under safety constraints.
Constrained Policy Optimization~\citep{achiam2017cpo} solves the CMDP
with a primal-dual approach inside a trust-region update.
\citet{stooke2020responsive} showed that the dual variable $\lambda$ can
be updated by a PID controller targeting the cost budget, yielding more
stable training than the pure proportional update of CPO; we build
directly on this PID-Lagrangian formulation. Recent infrastructure has
made the area accessible: Safety Gymnasium~\citep{ji2023safety} and
OmniSafe~\citep{ji2023omnisafe} provide unified benchmarks and
implementations of PPO-Lag and CPPO-PID, and
\citet{constraint_formulations_2024} surveyed the design space of
constraint formulations.

Safe RL has been applied to building energy management with mixed
success. \citet{safe_rl_buildings_mdpi_2025} train a barrier-certificate
agent that progressively reduces safety violations during training, and
\citet{safe_mbrl_hvac_2024} use a model-based RL system with uncertainty
gating for HVAC control. Neither work addresses the museum conservation
problem, the dynamic EN~15757 rate constraints, or the per-zone
collection heterogeneity studied here.

\subsection{Reinforcement learning for HVAC at large}

Deep RL for HVAC has progressed from early proof-of-concept studies on
single-zone office buildings~\citep{wang2017drlhvac,brandi2022online} to
recent benchmarks against ASHRAE Guideline~36 control sequences in
multi-zone buildings~\citep{drl_g36_2025}. Co-optimization of energy,
thermal comfort, and indoor air quality is the topic of
\citet{drl_co_optimize_2024}. Expert-guided warm-starting accelerates
training while preserving safety properties~\citep{efficient_assured_2025}.
\citet{boptest_rl_2025} validated an RL agent against real-building
measurements using the BOPTEST framework~\citep{boptest_mpc_2023}.

Multi-zone coordination is the next frontier.
\citet{drl_setpoint_multizone_2022} explored setpoint reset RL across
multiple zones; recent multi-agent RL work
\citep{marl_zone_alloc_2025,marl_electrochromic_2025,marl_demand_response_2025}
shows that centralized-training-with-decentralized-execution
architectures match or beat single-agent baselines while respecting
structural constraints. MPC reviews~\citep{drgona2020mpc_review} remain
the strongest classical baselines.

\subsection{Building digital twins and decarbonization}

Digital-twin (DT) technology is a maturing tool for building energy
management. Reviews \citep{dt_buildings_2025,dt_ieq_review_2025} report
that gaps between simulated and real building behavior often exceed
$30\%$ in occupant-driven variability. AI-enabled HVAC control is a
central element of net-zero strategies in 2025
roadmaps~\citep{ai_hvac_zeb_2025}, and DTs are increasingly paired with
LSTM and Transformer load forecasters~\citep{hybrid_lstm_transformer_load_2026,
lstm_transformer_power_2025}. For heritage applications specifically,
IoT-based microclimate monitoring has scaled from research prototypes
to deployable systems~\citep{iot_low_budget_museum_2022,
iot_microclimate_heritage_2024}.

\subsection{Building simulation testbeds}

Sinergym~\citep{manjavacas2024sinergym} provides a Gymnasium-style
wrapper around EnergyPlus~\citep{energyplus} and is the testbed we use
for high-fidelity validation. Alternative testbeds include
BOPTEST~\citep{boptest_mpc_2023} (Modelica-based) and the recent
BuildingGym~\citep{buildinggym_2025} toolbox. Our fast iteration
environment, a five-zone lumped-capacity thermal network, exposes the
same Gymnasium API as Sinergym, so the Lagrangian-PPO controller
transfers without code change to a calibrated EnergyPlus model.
Calibration follows the standard ASHRAE Guideline~14
metrics~\citep{ashrae_guideline14} (CV-RMSE${<}30\%$, NMBE${<}10\%$ for
hourly data).

% =====================  3. PROBLEM FORMULATION  =============================
\section{Problem Formulation}\label{sec:problem}

\subsection{Constrained Markov Decision Process}

We model the museum HVAC control problem as a discrete-time CMDP
$\langle\mathcal{S},\mathcal{A},P,r,c,\gamma,d\rangle$ with state space
$\mathcal{S}$, action space $\mathcal{A}$, transition kernel $P$, scalar
reward $r:\mathcal{S}\times\mathcal{A}\to\mathbb{R}$, scalar cost
$c:\mathcal{S}\times\mathcal{A}\to\mathbb{R}_{\ge 0}$, discount factor
$\gamma\in(0,1)$, and a cost budget $d\in\mathbb{R}_{\ge 0}$. A
stationary policy $\pi:\mathcal{S}\to\Delta(\mathcal{A})$ is feasible
when
\begin{equation}
\mathbb{E}_{\pi}\!\left[\sum_{t=0}^\infty \gamma^{\,t}\, c(s_t,a_t)\right]
\le d,
\label{eq:cmdp_constraint}
\end{equation}
and the constrained objective is
\begin{equation}
\max_{\pi}\;\;\mathbb{E}_{\pi}\!\left[\sum_{t=0}^\infty \gamma^{\,t}\, r(s_t,a_t)\right]
\;\text{ subject to \eqref{eq:cmdp_constraint}.}
\label{eq:cmdp_objective}
\end{equation}

\subsection{State, action, and reward}\label{sec:state_action_reward}

For a museum with $Z$ zones (we use $Z=5$ in our experiments) the state
at time $t$ aggregates four blocks:
\begin{equation}
s_t = \big[\,\mathbf{T}_t,\;\mathbf{H}_t,\;\boldsymbol{\theta}_t,\;
            \boldsymbol{u}_{t-1}\,\big]^{\top},
\label{eq:state}
\end{equation}
where $\mathbf{T}_t,\mathbf{H}_t\in\mathbb{R}^{Z}$ are zone air
temperature and relative humidity, $\boldsymbol{\theta}_t$ collects
outdoor weather, time-of-day, day-of-week, and instantaneous visitor
count, and $\boldsymbol{u}_{t-1}\in\mathbb{R}^{2Z}$ are the previous
control setpoints. The action $a_t\in[-1,1]^{2Z}$ is mapped affinely
onto per-zone temperature and humidity setpoints in
$[16,28]^\circ\mathrm{C}\times[30,70]\%\mathrm{RH}$.

The reward combines an energy term and a public-space comfort term:
\begin{equation}
r_t = -w_E\,\frac{E_t}{E_{\mathrm{norm}}} - w_C\,\max\big(0,|\mathrm{PMV}_t|-\eta\big),
\label{eq:reward}
\end{equation}
where $E_t$ is the per-step HVAC electrical energy (kWh),
$E_{\mathrm{norm}}$ a normalization constant, $\mathrm{PMV}_t$ the
Predicted Mean Vote in the public lobby, and $\eta=0.5$ a comfort
dead-band.

\subsection{Conservation cost}\label{sec:conservation_cost}

For zone $i$ with collection-determined envelope
$(T_{\min,i},T_{\max,i},\mathrm{RH}_{\min,i},\mathrm{RH}_{\max,i})$ and
EN~15757 rate caps $(\alpha_i,\beta_i)$, the per-step violation is

\begin{equation}
v^{T}_{i,t}  = \big[T_{\min,i}-T_{i,t}\big]_+ + \big[T_{i,t}-T_{\max,i}\big]_+,
\label{eq:vT}
\end{equation}
\begin{equation}
v^{\mathrm{RH}}_{i,t}= \big[\mathrm{RH}_{\min,i}-\mathrm{RH}_{i,t}\big]_+ +
                       \big[\mathrm{RH}_{i,t}-\mathrm{RH}_{\max,i}\big]_+,
\label{eq:vRH}
\end{equation}
\begin{equation}
v^{\dot T}_{i,t}     = \big[\,|\dot T_{i,t}|-\alpha_i\,\big]_+,
\label{eq:vTdot}
\end{equation}
\begin{equation}
v^{\dot{\mathrm{RH}}}_{i,t} = \big[\,|\dot{\mathrm{RH}}_{i,t}|-\beta_i\,\big]_+,
\label{eq:vRHdot}
\end{equation}
where $[\,\cdot\,]_+=\max(\,\cdot\,,0)$ and the rates $\dot T,
\dot{\mathrm{RH}}$ are computed by backward difference over the 15-min
control step. The scalar cost at time $t$ is the normalized sum
\begin{equation}
c_t = \sum_{i=1}^{Z}\Big(
   \tfrac{v^{T}_{i,t}}{\sigma_T}
 + \tfrac{v^{\mathrm{RH}}_{i,t}}{\sigma_{\mathrm{RH}}}
 + \tfrac{v^{\dot T}_{i,t}}{\sigma_{\dot T}}
 + \tfrac{v^{\dot{\mathrm{RH}}}_{i,t}}{\sigma_{\dot{\mathrm{RH}}}}\Big),
\label{eq:cost}
\end{equation}
with normalization scales $(\sigma_T,\sigma_{\mathrm{RH}},\sigma_{\dot T},
\sigma_{\dot{\mathrm{RH}}})=(1\,\mathrm{K},5\%,1\,\mathrm{K/h},2.5\%/\mathrm{h})$
chosen so that each component is $O(1)$ at a $1\sigma$ hard failure.

Table~\ref{tab:envelopes} lists the per-zone envelopes used in this
paper, derived from ASHRAE Handbook Chapter~24 class definitions
\citep{ashrae_handbook_ch24} and EN~15757~\citep{en15757}.

\begin{table}[t]
\centering
\caption{Per-zone conservation envelopes used in the simulator. ASHRAE
classes: A1 = strictest (organics on display), A2 = relaxed
(non-hygroscopic), --- = unconstrained.}
\label{tab:envelopes}
\small
\begin{tabular}{lccccccc}
\toprule
Zone & Class & $T_{\min}$ & $T_{\max}$ & $\mathrm{RH}_{\min}$
     & $\mathrm{RH}_{\max}$ & $\alpha$ & $\beta$\\
     &       & ($^\circ$C) & ($^\circ$C) & (\%) & (\%) & (K/h) & (\%/h)\\
\midrule
oil-painting gallery   & A1 & 19 & 24 & 45 & 55 & 2  & 5 \\
paper-works gallery    & A1 & 18 & 22 & 45 & 55 & 2  & 5 \\
bronze storage         & A2 & 15 & 25 & 30 & 40 & 3  & 10 \\
mixed gallery          & A2 & 18 & 24 & 40 & 60 & 3  & 10 \\
public lobby           & ---& 20 & 26 & 30 & 70 & 10 & 30 \\
\bottomrule
\end{tabular}
\end{table}

% =====================  4. METHODOLOGY  ====================================
\section{Methodology}\label{sec:method}

\begin{figure}[t]
\centering
\includegraphics[width=\linewidth]{figures/fig_framework.png}
\caption{System architecture of the Lagrangian-PPO museum HVAC
controller. The training pipeline is shown for $N$ parallel
environments. The simulator emits the state $s_t$ and the per-step
conservation cost $c_t$; the \texttt{CostRewardWrapper} fuses these
into a surrogate reward $\tilde r_t = r_t - \lambda c_t$; the PPO
actor-critic consumes the surrogate reward and emits the action $a_t$
that closes the inner loop. Once per rollout, the PID Lagrangian
controller updates $\lambda$ from the mean rollout cost
$\hat J_c$ and the cost budget $d$, and pushes the new value back to
every parallel \texttt{CostRewardWrapper} via the dashed feedback
edge (\texttt{set\_attr}). At deployment time, only the trained policy
is required.}
\label{fig:framework}
\end{figure}

Figure~\ref{fig:framework} summarizes the training and inference
pipeline. The simulator emits the state $s_t$ defined in
Eq.~\eqref{eq:state} together with the per-step conservation cost
$c_t$ from Eq.~\eqref{eq:cost}. A \texttt{CostRewardWrapper} subtracts
$\lambda c_t$ from the environment reward, and the resulting surrogate
reward in Eq.~\eqref{eq:lagrangian_reward} drives the PPO actor-critic
update. The Lagrange multiplier $\lambda$ is updated once per rollout
by the PID controller of Eq.~\eqref{eq:pid_lambda} and then propagated
to all parallel environments.

\subsection{Lagrangian PPO with a PID multiplier}

We solve the CMDP in Eq.~\eqref{eq:cmdp_objective} by Lagrangian
relaxation. For multiplier $\lambda\ge 0$, the surrogate per-step
reward is
\begin{equation}
\tilde r_t(\lambda) = r_t - \lambda\,c_t,
\label{eq:lagrangian_reward}
\end{equation}
and the saddle-point objective is
\begin{equation}
\min_{\lambda\ge 0}\;\max_{\pi}\;
\mathbb{E}_\pi\!\left[\sum_{t=0}^{\infty}\gamma^{\,t}\,
                 \big(r_t - \lambda\,(c_t - (1-\gamma)\,d)\big)\right].
\label{eq:saddle}
\end{equation}
The inner maximization is performed by Proximal Policy Optimization
(PPO) on the surrogate reward \eqref{eq:lagrangian_reward}, with the
standard Generalized Advantage Estimator~\citep{schulman2015gae}
($\gamma=0.99$, $\lambda_{\mathrm{GAE}}=0.95$), as implemented in
Stable-Baselines3~\citep{raffin2021sb3}. The outer minimization over $\lambda$ is
performed by a PID controller on the empirical violation
$\hat J_c - d$, where $\hat J_c$ is the mean per-step cost over the
most recent rollout~\citep{stooke2020responsive}:
\begin{equation}
\lambda \leftarrow \mathrm{clip}\Big(
   \lambda + K_P(\hat J_c - d) + K_I\!\!\sum_{k}(\hat J_{c,k}-d)
   + K_D\Delta(\hat J_c-d),\;\,
   [0,\lambda_{\max}]\Big).
\label{eq:pid_lambda}
\end{equation}
We use $K_P=0.05$, $K_I=5\times 10^{-4}$, $K_D=0$, and
$\lambda_{\max}=100$. Algorithm~\ref{alg:lpppo} summarizes the full
training loop.

\begin{algorithm}[t]
\caption{Lagrangian PPO with PID multiplier and subprocess-safe
$\lambda$ propagation.}
\label{alg:lpppo}
\begin{algorithmic}[1]
\State Initialize policy $\pi_\theta$, value $V_\phi$, multiplier
       $\lambda \leftarrow \lambda_0$, integrator $I\leftarrow 0$.
\State Create $N$ parallel envs $\{E_n\}_{n=1}^N$ wrapped with
       \texttt{CostRewardWrapper}, each holding an internal
       attribute \texttt{\_lambda\_scalar} initialized to $\lambda_0$.
\For{rollout $k=1,2,\dots$}
  \State Collect a rollout of length $T$ in each $E_n$ using the
         current policy; each step emits
         $\tilde r_t = r_t - E_n.\texttt{\_lambda\_scalar}\cdot c_t$
         and stores $c_t$ in an in-process per-env cost buffer.
  \State Run a PPO update on the rollout buffer (clipped surrogate
         + GAE) for $K$ epochs.
  \State \textbf{(End-of-rollout callback)} Pull
         $\{c_t\}$ buffers from all $N$ envs via
         \texttt{VecEnv.env\_method("pop\_cost\_buffer")}, compute
         $\hat J_c \leftarrow \mathrm{mean}(c)$.
  \State Update $\lambda$ via Eq.~\eqref{eq:pid_lambda}.
  \State \textbf{Push} $\lambda$ to all parallel envs via
         \texttt{VecEnv.set\_attr("\_lambda\_scalar",~$\lambda$)}.
\EndFor
\end{algorithmic}
\end{algorithm}

\subsection{Implementation pitfall: cross-process $\lambda$ propagation}\label{sec:pitfall}

Subprocess-based parallel environments such as Stable-Baselines3's
\texttt{SubprocVecEnv} construct each env inside a separate Python
process. An env factory closure is pickled, sent to the worker, and
unpickled before instantiation. A Python mutable object such as a
\texttt{dict} captured in the closure is therefore copied at
construction time, and any subsequent mutation by the parent process
is invisible to the workers.

\paragraph{Failure mode.}
We confirmed this experimentally. A naive implementation that stores
the multiplier in a shared \texttt{dict} produces six policies that
are byte-identical across cost budgets $d\in\{0,0.5,1,2,5,10\}$,
because every worker sees only the initial $\lambda_0$. The Pareto
front collapses to a single point. Diagnosis required comparing the
state dictionaries of independently trained networks.

\paragraph{Fix.}
Store $\lambda$ as a scalar attribute on each wrapper instance, and
push updates from the rollout-end callback via
\texttt{VecEnv.set\_attr}, which calls
\texttt{setattr(env, name, value)} inside each worker process. This
adds one round-trip per rollout, which is negligible compared to
thousands of env steps, and yields correct cross-process
synchronization. Algorithm~\ref{alg:lpppo} reflects the corrected
implementation.

We document this pitfall in detail because Safe RL building-energy
papers routinely report Lagrangian formulations without specifying the
vectorization strategy. A change between \texttt{DummyVecEnv}
(single-process; correct by default) and \texttt{SubprocVecEnv}
(multi-process; requires explicit synchronization) can silently
invalidate the constraint.

% =====================  5. EXPERIMENTAL SETUP  ==============================
\section{Experimental Setup}\label{sec:setup}

\subsection{Simulator}\label{sec:simulator}

For fast experimentation we implement a five-zone lumped-capacity
museum simulator. Each zone has independent thermal mass and a
first-order RC response to (i) the HVAC setpoint with time constant
$\tau_{\text{HVAC}}=0.5\,\mathrm{h}$, (ii) the outdoor envelope with
$\tau_{\text{env}}=8\,\mathrm{h}$, and (iii) occupant heat gains scaled
by per-zone occupancy share. Humidity follows an analogous dynamic with
a faster setpoint response constant. Outdoor weather is a synthetic
seasonal sinusoid superposed with a diurnal cycle and Gaussian noise.
The control timestep is $15\,\mathrm{min}$, and an episode covers one
calendar year (35\,040 steps).

The simulator observation and action spaces, conservation envelopes,
and reward composition are those of
Section~\ref{sec:state_action_reward}. The same wrapper API supports a
Sinergym/EnergyPlus backend with the zone-to-collection mapping in
Table~\ref{tab:envelopes}, so the pipeline can be re-run on the
calibrated high-fidelity model without algorithm changes. We report
fast-iteration results in this paper and discuss the EnergyPlus
transition in Section~\ref{sec:discussion}.

\subsection{Visitor flow}

Visitor counts are synthesized from a weekly schedule (Mondays closed,
weekend boost) combined with a smooth daily bump between 09:00 and
17:00, plus Gaussian noise. The per-step occupancy enters the thermal
balance as a sensible and latent heat gain proportional to per-zone
share.

\subsection{Baselines}

We compare five controllers that span the spectrum from naive rule-based
control to safe RL:

\begin{description}[leftmargin=1.2em,style=nextline,nosep]
  \item[B1: Rule-fixed.] Constant setpoint $22^\circ\mathrm{C}/50\%\mathrm{RH}$ in all zones, all hours.
  \item[B2: Schedule.] Setpoints relaxed during closed hours (Mondays
        and 18:00--08:00) to wider per-zone deadbands.
  \item[B3: Vanilla-PPO.] PPO trained on energy and comfort only; no
        conservation term in the reward.
  \item[B4: Soft-PPO.] PPO with a fixed-weight soft penalty
        $r-w\,c$, $w=5.0$.
  \item[B5: Lagrangian-PPO (ours).] Algorithm~\ref{alg:lpppo}.
\end{description}

All RL baselines use the same MLP policy with two hidden layers of
256 units, learning rate $3\times 10^{-4}$, GAE
$\lambda_{\mathrm{GAE}}=0.95$, discount $\gamma=0.99$, clip range
$0.2$, and batch size $256$. Each of B3--B5 is trained for
$5\times 10^{5}$ steps under 16 parallel envs using
Stable-Baselines3~\citep{raffin2021sb3}.

\subsection{Pareto sweep}

To trace the energy--conservation trade-off, we additionally train
Lagrangian-PPO for $2\times 10^{5}$ steps at six cost budgets
$d\in\{0,0.5,1,2,5,10\}$ and report each policy's deterministic
year-long evaluation.

\subsection{Evaluation metrics}\label{sec:metrics}

Each trained policy is evaluated on a single deterministic year-long
rollout from a fixed seed. The reported metrics are annual HVAC energy
(kWh and MWh), the year-total conservation cost from
Eq.~\eqref{eq:cost}, the number of violation steps per year (steps
with $c_t>0$, each step is 15\,min), the per-zone decomposition of
violation steps, the mean PMV in the public lobby, and the number of
comfort-violation steps.

% =====================  6. RESULTS  =========================================
\section{Results}\label{sec:results}

\subsection{Headline comparison}

Table~\ref{tab:headline} summarizes annual energy and conservation
metrics across all controllers, aggregated over three seeds for the
trained controllers. The Lagrangian-PPO controller trained with the
strictest budget ($d=0$) reduces year-total conservation cost by
\SI{85}{\percent} relative to the fixed-setpoint baseline B1 at the
cost of a \SI{1.1}{\percent} energy increase, and reduces violation
steps from 35\,040 (the entire year) to 20\,738, or about
\SI{59}{\percent} of the year.

Both unconstrained variants (B3 vanilla, B4 soft) achieve nominally
lower annual energy than B1 but inflict $3$--$6\times$ higher
conservation cost: the unconstrained reward signal drives the policy
toward setpoints that minimize HVAC work without protecting
collections.

The $500\mathrm{k}$-step Lagrangian variant (\texttt{lagr}) actually
underperforms the $200\mathrm{k}$-step strict-budget sweep on
conservation cost (\num{157}~000 vs.\ \num{11.7}~000). The cause is
$\lambda$ saturation: once $\lambda=\lambda_{\max}$, the policy is
pushed into a corner of action space that minimizes integrated cost at
the expense of robustness, and additional training amplifies the
effect. We return to this in Section~\ref{sec:discussion}.

\begin{table}[t]
\centering
\caption{Year-long evaluation of all controllers (deterministic
rollouts). Mean$\pm$SEM across three seeds where applicable.
Conservation cost is the year-total of Eq.~\eqref{eq:cost}.
Violation steps is the number of 15-min steps with non-zero cost,
out of 35\,040 in a year.}
\label{tab:headline}
\small
\setlength{\tabcolsep}{4pt}
\begin{tabular}{lcccc}
\toprule
Controller & Energy (MWh) & Cons.\ cost & Viol.\ steps & Mean PMV \\
\midrule
B1 Rule-fixed       & 364.8                       & 78{,}380   & 35{,}040 & $-0.32$ \\
B2 Schedule         & 350.7                       & 75{,}833   & 35{,}040 & $-1.54$ \\
B3 Vanilla-PPO      & $361.1\pm 3.1$              & 348{,}091  & 35{,}040 & $-0.10$ \\
B4 Soft-PPO ($w=5$) & $380.5\pm 0.3$              & 496{,}192  & 35{,}040 & $-0.25$ \\
B5 Lagr.-PPO 500k   & $372.4\pm 1.6$              & 157{,}033  & 35{,}040 & $\phantom{-}$varies \\
\midrule
\multicolumn{5}{l}{\emph{Pareto sweep, single seed, 200k steps:}}\\
$d=0.0$  & 368.8 & \textbf{11{,}677} & 20{,}738 & $0.65$ \\
$d=0.5$  & 368.9 & 12{,}197          & 21{,}161 & $0.66$ \\
$d=1.0$  & 369.0 & 12{,}807          & 21{,}625 & $0.66$ \\
$d=2.0$  & 369.1 & 13{,}911          & 22{,}475 & $0.67$ \\
$d=5.0$  & 369.4 & 18{,}450          & 26{,}789 & $0.70$ \\
$d=10.0$ & 370.2 & 28{,}684          & 35{,}003 & $0.75$ \\
\bottomrule
\end{tabular}
\end{table}

\subsection{Pareto front}

Figure~\ref{fig:hero}(c) traces the Pareto front of Lagrangian-PPO
solutions as the cost budget $d$ ranges from $0$ to $10$. The front
is strongly anisotropic: annual energy varies by only
$1.4\,\mathrm{MWh}$ (\SI{0.4}{\percent}) across the six budgets, while
the conservation cost varies by a factor of $2.5$. This anisotropy is
a direct consequence of the rate clauses. Most of the energy required
to hold the static envelope, Eqs.~\eqref{eq:vT}--\eqref{eq:vRH}, is
already paid for by any controller that keeps $T$ inside the envelope
at all, whereas the rate clauses in
Eqs.~\eqref{eq:vTdot}--\eqref{eq:vRHdot} can be satisfied by smoother
trajectories with negligible additional energy. In other words, how
strictly to enforce conservation has order-of-magnitude larger
conservation impact than energy impact under the EN~15757 formulation.

\begin{figure}[t]
\centering
\includegraphics[width=\linewidth]{figures/fig1_hero.pdf}
\caption{Year-long evaluation of all controllers on the five-zone
museum simulator. (a) Annual HVAC energy. (b) Number of 15-min steps
with at least one zone violating its conservation envelope, out of
35\,040 in a year. (c) Pareto front of Lagrangian-PPO solutions
($d\in\{0,0.5,1,2,5,10\}$, color-coded). The front is strongly
anisotropic: annual energy varies by less than $\SI{0.5}{\percent}$
across the budget range while conservation cost spans a factor of
$2.5$. (d) Per-zone breakdown of violation steps. The strict-budget
Lagrangian controllers confine residual violations to a single
constraint-coupled zone (\texttt{paper\_works\_gallery}), whereas
unconstrained or saturated controllers spread violations across all
five zones.}
\label{fig:hero}
\end{figure}

\subsection{Typical-week trajectories}

Figure~\ref{fig:typical_week} compares the per-step conservation cost
and HVAC energy of the Rule-fixed baseline (B1) and the strict-budget
Lagrangian-PPO controller across a representative week. The
Lagrangian controller absorbs the daily and weekly disturbances with
smoother control trajectories; the cost trace stays close to zero
during steady-state daytime operation, whereas the rule-fixed
controller incurs a persistent baseline cost from the unavoidable
oil--bronze envelope conflict.

\begin{figure}[t]
\centering
\includegraphics[width=0.92\linewidth]{figures/fig2_typical_week.pdf}
\caption{Per-step conservation cost (top) and HVAC energy (bottom) for
the Rule-fixed baseline (gray) and the strict-budget Lagrangian-PPO
controller (red), over a representative one-week window.}
\label{fig:typical_week}
\end{figure}

\subsection{Per-zone localization}

The per-zone decomposition in Figure~\ref{fig:hero}(d) is the
clearest qualitative result. The constrained controllers
(sweep $d\in\{0,\dots,5\}$) restrict essentially all residual
violations to the \texttt{paper\_works\_gallery} zone, which is the
only zone whose temperature envelope ($18^\circ\mathrm{C}$ to
$22^\circ\mathrm{C}$) overlaps the public-lobby comfort band and the
bronze-storage humidity band in a way that the shared HVAC supply
cannot simultaneously satisfy. The unconstrained variants
(\texttt{vanilla}, \texttt{soft}) and the saturated $500$k-step
Lagrangian variant (\texttt{lagr}) instead distribute violations
across all five zones, including \texttt{oil\_painting\_gallery} and
\texttt{mixed\_gallery}, which would never violate under any
reasonable controller.

This localization property has practical importance. When residual
violations are isolated to one zone, conservation engineering can
focus on that zone, for instance by adding a local dehumidifier or a
revised duct layout, without re-tuning the controller for the entire
building. When violations are distributed, no single retrofit can
address them.

\subsection{$\lambda$ saturation}\label{sec:lambda_saturation}

A surprising finding is the deteriorating performance of the
$500$k-step Lagrangian runs (Table~\ref{tab:headline}, row
\texttt{lagr}) relative to the $200$k-step sweep at $d=0$. Both target
the same cost budget; the only difference is training duration and
the resulting $\lambda$ trajectory. The $500$k-step run saturates at
$\lambda_{\max}=100$ after roughly $1\times 10^5$ steps, while the
$200$k-step run converges to $\lambda\approx 69$ without saturation.

Once $\lambda$ saturates, $\lambda c_t$ dominates the surrogate reward
in Eq.~\eqref{eq:lagrangian_reward}, and PPO is effectively trained to
minimize cost while ignoring the original reward. The policy is
pushed to an extreme corner of the action space where the conservation
constraint is approximately satisfied but the controller no longer
generalizes to off-distribution states. Deterministic evaluation on a
held-out year then traces trajectories that violate more zones than
the unsaturated policy. Figure~\ref{fig:training_curves} shows the
$\lambda$ trajectory and the mean per-rollout cost across the three
Lagrangian seeds: $\lambda$ rises monotonically and saturates after
approximately $1\times 10^5$ steps, after which the mean cost
ceases to decrease and exhibits a long plateau.

\begin{figure}[t]
\centering
\includegraphics[width=\linewidth]{figures/fig3_training_curves.pdf}
\caption{Training dynamics of the Lagrangian multiplier and mean cost.
(a) $\lambda$ trajectory: monotone rise toward $\lambda_{\max}=100$,
saturating around step $1\times 10^5$. (b) Mean rollout cost: the
controller drives cost down during early training, but the rate of
descent stalls once $\lambda$ saturates. Different curves correspond
to the three independent training seeds.}
\label{fig:training_curves}
\end{figure}

A practical remedy is to set $\lambda_{\max}$ conservatively. We
suggest $\lambda_{\max}=30$ as a starting point for the rate-cap
formulation. We also tested a constraint-decomposed Lagrangian with one
multiplier per envelope clause, expecting it to relieve the saturation,
but it performed markedly worse than the scalar formulation: the
independent per-clause PID loops fail to coordinate, the rate-clause
multipliers run away while the static-clause multipliers stay small, and
the resulting policy violates the static envelopes across all zones. The
scalar cost, by balancing the four clause families through a single
multiplier, is the more robust choice here.

\subsection{Validation on Sinergym + EnergyPlus}\label{sec:sinergym_validation}

To test whether the mock-env conclusions transfer to an
engineering-grade model, we built a high-fidelity testbed on
Sinergym's~\citep{manjavacas2024sinergym} canonical five-zone
EnergyPlus building (NY-JFK TMY3 weather, full calendar year, 15-min
step).

\paragraph{Single-zone diagnosis.}
The stock IDF drives all five zones from one shared thermostat schedule
and includes no active humidification. Under these conditions every
controller, rule-based and learned alike, violates the conservation
envelope at essentially every step, and a cost-budget sweep collapses
to indistinguishable trajectories. A clause-level decomposition of the
cost explains why: about $95\%$ of the conservation cost comes from the
static humidity term. Without a humidifier, indoor relative humidity
averages $26\%$ in winter, far below the ASHRAE Class~A1 band
($45$--$55\%$), so the humidity envelope is physically unsatisfiable by
any thermostat policy. This corrects a natural but mistaken reading of
the single-zone result as a reinforcement-learning failure: the
binding obstacle is missing equipment, not the optimizer.

\paragraph{Multi-zone humidified testbed.}
We therefore modified the building to reflect how museum HVAC is
actually provisioned. We replaced the shared thermostat with five
independent per-zone dual-setpoint thermostats (a 10-dimensional
setpoint action), and added a central steam humidifier driven by a
humidistat so the building can hold a humidity setpoint. With
humidification the indoor relative humidity rises to a $51\%$ annual
mean, inside the Class~A1 band, which makes the conservation
constraints physically feasible. The five EnergyPlus zones are mapped
onto the five collection envelopes of Table~\ref{tab:envelopes}.
Table~\ref{tab:sinergym} reports the year-long evaluation of the
fixed-setpoint baseline and the Lagrangian-PPO controller on this
model.

\begin{table}[t]
\centering
\caption{Multi-zone humidified Sinergym/EnergyPlus validation (five
per-zone thermostats plus a central steam humidifier, NY-JFK TMY3
weather). Lagrangian-PPO is the mean over three seeds; the three seeds
converge to the same policy. Per-zone counts are violation steps out of
35\,040.}
\label{tab:sinergym}
\small
\setlength{\tabcolsep}{4pt}
\begin{tabular}{lcc}
\toprule
& Rule-fixed & Lagr.-PPO \\
\midrule
Energy (MWh)            & 36.7    & 20.7 \\
Conservation cost       & 144{,}643 & 265{,}338 \\
Violation steps         & 35{,}040  & 35{,}040 \\
\midrule
oil-painting gallery    & 7{,}554   & 18{,}674 \\
paper-works gallery     & 22{,}239  & 30{,}793 \\
bronze storage          & 35{,}033  & 35{,}018 \\
mixed gallery           & 4{,}292   & 13{,}144 \\
public lobby            & 504       & 10{,}128 \\
\bottomrule
\end{tabular}
\end{table}

\paragraph{Localization holds on the high-fidelity model.}
The per-zone localization seen on the mock simulator reappears here.
Under the fixed-setpoint baseline, residual violations concentrate in
the bronze-storage zone ($35{,}033$ of $35{,}040$ steps) while the
public lobby violates only $504$ steps. The central humidifier holds
the supply air near $50\%$ relative humidity, which satisfies the oil-
and paper-gallery bands ($45$--$55\%$) but is structurally too humid
for the bronze envelope ($30$--$40\%$). The oil/bronze humidity
incompatibility that motivates the per-zone formulation is therefore a
real property of the building physics, not an artifact of the lumped
model.

\paragraph{Reinforcement learning does not beat the fixed setpoint here.}
Contrary to the mock-env results, the Lagrangian-PPO controller does
not improve on the fixed-setpoint baseline on this model. It lowers
annual energy ($20.7$ vs.\ $36.7$\,MWh) but raises the year-total
conservation cost ($265{,}338$ vs.\ $144{,}643$) and incurs more
violation steps in every collection zone. The Lagrange multiplier
saturates at $\lambda_{\max}$ and the three seeds converge to the same
energy-leaning policy. On the EnergyPlus model the energy gradient
dominates and the bounded-action PPO settles into an energy-saving
region that the cost penalty does not pull back, a behavior we did not
observe on the better-conditioned mock environment and that a
tanh-squashed SAC variant did not resolve either.

\paragraph{Interpretation.}
The high-fidelity study tempers the mock-env conclusions. The CMDP
formalization, the per-zone localization insight, and the qualitative
energy/conservation trade-off survive the move to EnergyPlus, but the
stronger claim that constrained RL reduces conservation cost relative
to a fixed setpoint is supported only on the mock simulator. On the
EnergyPlus model the learned controllers trade conservation for energy
and are outperformed by a well-chosen fixed setpoint. Closing this gap
likely requires reward normalization and multiplier scheduling so the
constraint signal survives the strong energy gradient, which we leave
to future work.

% =====================  7. DISCUSSION  =====================================
\section{Discussion}\label{sec:discussion}

\subsection{Limitations}

\paragraph{Simulator fidelity.}
The five-zone lumped-capacity model used here is a fast-iteration
testbed, not an engineering-grade representation. Real museum HVAC
includes coupled air-handling units, return-air mixing, fresh-air
penalties, and structural infiltration that the lumped model omits.
Our wrapper API exposes the same observation and action interface as
Sinergym~\citep{manjavacas2024sinergym}, so substituting the
EnergyPlus backend is a single-line change in the env factory. We
report only lumped-model results here for reproducibility and rapid
iteration, and we plan to re-run the full pipeline on a calibrated
EnergyPlus model in follow-up work.

\paragraph{Visitor model.}
The synthetic visitor flow captures weekly and daily seasonality but
omits visitor-mass effects on humidity gain and the slow drift in
attendance with weather and special exhibitions. A real-data
deployment would replace this module with an LSTM-based or
Transformer-based forecaster trained on the museum's ticketing
records.

\paragraph{Constraint shaping.}
We use a scalar cost in Eq.~\eqref{eq:cost} that sums four normalized
components. The multi-constraint Safe RL literature would instead treat
each clause as an independent constraint with its own $\lambda$. We
implemented and tested this constraint-decomposed variant
(Section~\ref{sec:lambda_saturation}) but found it less robust than the
scalar cost: without coordination across the per-clause multipliers,
one clause dominates and the others are starved. A principled
decomposition would need coupled or budget-shared multiplier updates,
which we leave to future work.

\subsection{When does Safe RL beat classical multi-objective optimization?}

The closest classical baselines~\citep{ascione2018museum,schito2020museum}
solve a multi-objective optimization at design time and produce HVAC
setpoint schedules that are then applied open-loop. Our results clarify
the conditions under which closed-loop RL adds value. When loads
(visitors, weather) vary on the same time scale as the HVAC response
(hours), open-loop schedules are forced to be conservative; closed-loop
control can ride the disturbances. When the building has heterogeneous
collections, the per-zone reward shaping that RL allows (via
independent $\lambda$ per envelope) is harder to express in a single
Pareto-front problem. When the constraint is dynamic, as with the
EN~15757 rates, encoding it as a CMDP cost is more direct than
treating it as a post-hoc filter on a Pareto-optimal schedule.

In settings where these conditions do not apply, for example a small
single-zone storage room with a stable load, classical multi-objective
optimization may remain the preferred tool. Safe RL should be regarded
as a complement to, not a replacement for, the classical pipeline.

\subsection{Why reinforcement learning underperforms on EnergyPlus}

The gap between the mock and EnergyPlus results
(Section~\ref{sec:sinergym_validation}) is informative. On the mock
simulator the energy and conservation terms sit on comparable scales,
the dynamics are smooth, and the Lagrangian multiplier steers the
policy onto the constraint boundary. On the EnergyPlus model the
energy gradient is much stronger relative to the conservation cost,
and the bounded-action PPO settles into an energy-saving region early
in training; once the policy entropy has collapsed, raising $\lambda$
to its ceiling no longer moves it, and a tanh-squashed SAC variant
that preserves exploration does not recover the constraint either.
Two observations follow. First, the methodological pieces that do not
depend on the optimizer transfer cleanly: the EN~15757 CMDP cost, the
per-zone differentiation, and the localization of residual violations
to the humidity-incompatible bronze zone all hold on the high-fidelity
model. Second, making constrained RL actually beat a tuned fixed
setpoint on EnergyPlus appears to require constraint-aware conditioning
of the reward, for example per-term normalization of the energy and
cost signals and a multiplier schedule that ramps before the policy
entropy decays. We regard this as the central open problem for
deploying the method beyond the mock testbed.

\subsection{A reusable high-fidelity museum testbed}

A by-product of this study is a high-fidelity museum HVAC testbed that
we release with the code. It augments Sinergym's five-zone EnergyPlus
building with five independent per-zone dual-setpoint thermostats and a
central steam humidifier, and maps the five zones onto the collection
envelopes of Table~\ref{tab:envelopes}. The same
\texttt{CostRewardWrapper} and Lagrangian callback drive both the mock
and the EnergyPlus backends without change. The remaining step toward a
deployable model is calibration to a specific museum's monthly energy
records following ASHRAE Guideline~14~\citep{ashrae_guideline14}
(CV-RMSE${<}30\%$, NMBE${<}10\%$), which requires building-specific
metering data beyond the scope of this paper. We provide the testbed so
that the energy/conservation conditioning problem identified above can
be studied on a realistic model rather than only on a lumped one.

% =====================  8. CONCLUSION  =====================================
\section{Conclusion}\label{sec:conclusion}

We formalized the museum HVAC control problem as a Constrained
Markov Decision Process that encodes both ASHRAE-style static
temperature and humidity envelopes and EN~15757 dynamic rate clauses
as a single scalar cost. We solved it with a Lagrangian-PPO
controller whose multiplier is driven by a PID controller, and we
documented an implementation issue ($\lambda$ propagation across
parallel subprocess envs) that silently invalidates the constraint
in a naive implementation.

On a five-zone differentiated museum simulator we obtained a clean
three-objective Pareto front whose anisotropy, in which conservation
cost varies by a factor of $2.5$ while energy varies by less than
$0.5\%$ across budgets, gives operators a principled tool for trading
off collection safety against energy. At the strictest budget the
controller reduces year-total conservation cost by $85\%$ relative to a
fixed-setpoint baseline at a $1\%$ energy overhead, and a per-zone
analysis shows that residual violations stay confined to a single
constraint-coupled zone.

We then stress-tested these conclusions on a high-fidelity EnergyPlus
model, which we extended with per-zone thermostats and active
humidification so that the conservation constraints become physically
feasible. The localization result transfers: residual violations
concentrate in the humidity-incompatible bronze-storage zone. The
performance result does not: on EnergyPlus the constrained controller
trades conservation for energy and is outperformed by a tuned fixed
setpoint, because the strong energy gradient overwhelms the constraint
signal and the multiplier saturates. We report this honestly. The
formalization, the localization insight, and the released testbed stand
on their own, while making constrained RL beat a fixed setpoint on the
high-fidelity model remains open and, we argue, hinges on
constraint-aware conditioning of the reward rather than on the
constraint encoding itself.

Future work will (i) calibrate the EnergyPlus testbed to a real
museum's metering records, (ii) study reward normalization and
multiplier scheduling to close the mock-to-EnergyPlus gap, and (iii)
incorporate real visitor forecasts via time-series transformers.

% =====================  ACKNOWLEDGMENTS  ====================================
